\chapter{PENDAHULUAN}

\section{Latar Belakang}
Bagian ini menjadi pembuka seluruh dokumen. Gunakan paragraf pertama untuk menjelaskan konteks umum topik yang Anda angkat — mengapa topik ini relevan, di bidang apa ia berada, dan apa keadaan atau permasalahan yang melatarbelakanginya. Jangan langsung masuk ke detail teknis; mulailah dari gambaran besar lalu perkecil ke masalah spesifik yang ingin Anda selesaikan.

Paragraf kedua dan seterusnya sebaiknya merujuk sumber-sumber yang mendukung argumen Anda. Setiap kali Anda menyebut sebuah fakta, teori, atau hasil penelitian orang lain, sertakan sitasi agar pembaca bisa melacak asal informasi tersebut. Template ini mendukung dua gaya sitasi: sitasi naratif yang menyatu dalam kalimat, dan sitasi dalam kurung yang berdiri sendiri. Kedua gaya tersebut mengambil data dari file bibliography yang sama, sehingga Anda cukup menambahkan entri sekali lalu memanggilnya berkali-kali di berbagai bagian dokumen.

Hindari menumpuk banyak sitasi dalam satu kalimat. Satu atau dua sumber per kalimat sudah cukup; lebih dari itu akan membuat paragraf terasa sesak dan menurunkan keterbacaan. Sebagai contoh, sitasi naratif menghasilkan Silberschatz et al.\ (2020) \textcite{silberschatz2020}, sedangkan sitasi dalam kurung menghasilkan (Masri and Suhartini, 2021) \parencite{masri2021}. Kedua gaya mengambil data dari file yang sama, sehingga satu entri bisa dipanggil berkali-kali di berbagai bagian dokumen.

\section{Rumusan Masalah}
Rumusan masalah adalah inti dari dokumen ini. Tuliskan dalam bentuk pertanyaan yang jelas dan terukur, atau dalam bentuk pernyataan masalah yang spesifik. Setiap rumusan masalah akan dijawab kembali di bagian hasil dan penutup, jadi pastikan jumlahnya proporsional dengan kedalaman pekerjaan yang Anda lakukan.

Gunakan penomoran berurutan agar pembaca mudah melacaknya. Berikut contoh formatnya:

\begin{enumerate}
    \item Bagaimana merancang sistem X yang dapat memenuhi kebutuhan Y?
    \item Bagaimana mengevaluasi kinerja sistem X setelah diimplementasikan?
    \item Apakah hasil evaluasi menunjukkan peningkatan dibandingkan kondisi awal?
\end{enumerate}

Dua hingga empat rumusan masalah biasanya sudah memadai untuk laporan standar; lebih dari itu pertimbangkan untuk memecah dokumennya.

\section{Tujuan}
Tujuan ditulis sebagai jawaban langsung dari rumusan masalah. Jika rumusan masalahnya berbentuk pertanyaan, maka tujuannya adalah menyatakan apa yang ingin dicapai untuk menjawab pertanyaan tersebut. Pastikan setiap tujuan dapat diukur atau diverifikasi melalui hasil yang Anda sajikan di bab-bab berikutnya.

Tulis setiap tujuan dalam bentuk kalimat yang menyatakan apa yang akan dilakukan, tanpa penomoran khusus — cukup urutkan secara logis mengikuti urutan rumusan masalah.

\section{Manfaat}
Bagian ini menjelaskan untuk siapa dan dalam konteks apa dokumen ini berguna. Pisahkan manfaat teoretis (kontribusi terhadap pemahaman atau literatur) dari manfaat praktis (dapat langsung diterapkan oleh pembaca atau pihak terkait).

Tidak perlu berlebihan; dua hingga tiga manfaat sudah cukup. Hindari manfaat yang terlalu umum tanpa konteks spesifik.
